\section{Gaps, Charges, and Translation Sectors}
\label{sec:gaps-charges-sectors}

\subsection{Cyclic gaps and excess charge}
\label{subsec:cyclic-gaps-excess-charge}

Let \(Q\in\{\pm1\}^n\) be a cyclic quadrilateral-flux word that admits a
cyclic \(\tau\)-lift.  Suppose first that its positive set
\[
D(Q)=\{i\in\mathbb Z/n\mathbb Z:Q_i=+1\}
\]
is nonempty.  Choose integer representatives
\[
x_1<x_2<\cdots<x_d<x_1+n
\]
for its \(d\) elements and put \(x_{d+1}=x_1+n\).  The positive cyclic gaps
and their excess charges are
\[
g_j=x_{j+1}-x_j,\qquad q_j=g_j-4
\qquad(1\leq j\leq d).
\]
The reference value is four; positive or negative \(q_j\) measures the
departure of one gap from that spacing.

\begin{proposition}[Gap and charge identities]\label{prop:gap-charge}
The gap list satisfies
\[
\sum_{j=1}^d g_j=n,\qquad
\sum_{j=1}^d q_j=n-4d.
\]
If \(n\) is even and \(Q\) has a cyclic \(\tau\)-lift, then \(d\) is even
and
\[
\sum_{j=1}^d q_j\equiv n\pmod8.
\]
\end{proposition}

\begin{proof}
The half-open arcs from \(x_j\) to \(x_{j+1}\) partition one traversal of
the \(n\)-cycle, so their lengths sum to \(n\).  Subtracting four for each
of the \(d\) arcs gives the second identity.

The word has \(d\) positive and \(n-d\) negative entries.  Hence
\[
\prod_{i=0}^{n-1}Q_i=(-1)^{n-d}.
\]
A cyclic lift exists only when this product is \(1\).  Thus \(n-d\) is
even; since \(n\) is even, \(d\) is even.  Write \(d=2h\).  The second
identity becomes
\[
\sum_{j=1}^d q_j=n-8h,
\]
which is the asserted congruence.
\end{proof}

The charge law depends only on \(Q\), while \(\alpha\) is an independent
step-one cycle sign. The defect-free word \(D(Q)=\varnothing\), which has no
cyclic gap list, is handled outside Proposition~\ref{prop:gap-charge}.

\subsection{Translation sectors}
\label{subsec:translation-sectors}

Recall the four reference translates \(B_s\), \(s\in\mathbb Z/4\mathbb Z\),
defined by
\[
(B_s)_i=+1\quad\Longleftrightarrow\quad i\equiv s\pmod4.
\]
They are sectors of the \(Q\)-word.  Each has two period-eight
\(\tau\)-lifts, but the sector label itself depends only on the positions of
the positive \(Q\)-sites.

\begin{proposition}[Oriented gap shift]\label{prop:sector-shift}
Suppose an oriented gap of length \(g\) leaves the reference sector \(B_s\)
at a positive site and enters a reference sector on its right.  For its
charge \(q=g-4\), the right sector is
\[
B_{s+\sigma_{\mathrm{sec}}(q)},\qquad
\sigma_{\mathrm{sec}}(q)=q\pmod4.
\]
This law is independent of the \(\tau\)-lift and Hamilton holonomy.
\end{proposition}

\begin{proof}
Let the left endpoint be \(x\equiv s\pmod4\).  The next positive site is
\(x+g\), so the positive lattice continuing to the right is
\[
x+g+4\mathbb Z.
\]
It is therefore \(B_{s+g}\). Since \(g=q+4\), the residues of \(g\) and
\(q\) agree modulo four. Because the proof uses only endpoint positions, the
sector shift is independent of the lift sign and \(\alpha\).
\end{proof}

\subsection{Composition and legal cyclic closure}
\label{subsec:composition-legal-cyclic-closure}

For consecutive oriented gaps with charges \(q_1,\ldots,q_k\), repeated use
of Proposition~\ref{prop:sector-shift} gives
\[
B_s\longmapsto B_{s+q_1}\longmapsto\cdots\longmapsto
B_{s+q_1+\cdots+q_k},
\]
with labels reduced modulo four.  Thus sector shifts add:
\[
\sigma_{\mathrm{sec}}(q_1+\cdots+q_k)
\equiv\sum_{j=1}^k\sigma_{\mathrm{sec}}(q_j)\pmod4.
\]
On a cyclic gap list, \(\sum_jg_j=n\), so the final reference lattice is
the translate \(B_{s+n}\); equivalently,
\[
\sum_jq_j\equiv n\pmod4.
\]
This compatibility is weaker than the mod-eight law in Proposition
\ref{prop:gap-charge}, because it does not contain the parity information
needed for a cyclic \(\tau\)-lift.

The modulo-eight charge closure and the modulo-four translation-sector law
encode different constraints and will be used separately.  The first is a
global condition on a liftable ring; the second records the local reference
translate across an interface.

For later use, a gap of length six has charge \(q=6-4=2\) and sector shift
two modulo four.  One, two, or three such gaps contribute total charges
\(2,4,6\), respectively, and sector shifts \(2,0,2\) modulo four.  If all
other gaps have the reference length four, these are precisely the charge
contributions required for ring orders congruent to \(2,4,6\) modulo eight.
No spectral comparison is asserted at this stage.
